%% ACM sigconf format — 11 pages target
\documentclass[sigconf,10pt]{acmart}

\usepackage{listings}
\usepackage{booktabs}
\let\Bbbk\relax
\usepackage{amsmath,amssymb}
\usepackage{xcolor}
\usepackage{microtype}
\usepackage{balance}

%% Tolerate slightly loose lines to avoid overfull \texttt identifiers
\tolerance=1500
\emergencystretch=1em

%% Lean 4 listing style
\lstdefinelanguage{Lean4}{
  keywords={theorem, def, lemma, private, inductive, structure, where, namespace, end,
            import, open, sorry, by, have, let, in, match, with, if, then, else,
            return, partial, fun, forall, exists, intro, exact, apply, refine,
            constructor, induction, cases, simp, omega, rw, ext, congr, tauto,
            decide, native_decide, aesop, rcases, obtain, use, suffices},
  sensitive=true,
  morecomment=[l]{--},
  morecomment=[n]{/-}{-/},
  morestring=[b]",
}
\lstset{
  language=Lean4,
  basicstyle=\ttfamily\scriptsize,
  keywordstyle=\bfseries\color{blue!70!black},
  commentstyle=\itshape\color{gray!80},
  stringstyle=\color{green!50!black},
  breaklines=true,
  frame=single,
  framerule=0.4pt,
  xleftmargin=4pt,
  xrightmargin=4pt,
  aboveskip=4pt,
  belowskip=4pt,
  literate=
    {≠}{{$\neq$}}1
    {→}{{$\to$}}1
    {∃}{{$\exists$}}1
    {∈}{{$\in$}}1
    {∧}{{$\land$}}1
    {∀}{{$\forall$}}1
    {≥}{{$\geq$}}1
    {₁}{{$_1$}}1
    {₂}{{$_2$}}1,
}

\lstdefinelanguage{Rust2}{
  keywords={pub, fn, let, mut, if, else, return, match, for, in, while,
            struct, enum, impl, use, mod, type, trait, where, self, Self,
            true, false, Some, None, Ok, Err},
  sensitive=true,
  morecomment=[l]{//},
  morecomment=[n]{/*}{*/},
  morestring=[b]",
}

\begin{document}

\title{Truly Automated Formal Specification and Verification with Near-Zero Human Labour}

\author{Don Syme}
\affiliation{\institution{GitHub Next}\country{United Kingdom}}
\email{dsyme@github.com}

\author{LeanSquad}
\affiliation{\institution{AI Automation Workflow (GitHub Agentic Workflows)}\country{}}
\email{lean-squad@automated}

\author{Copilot CLI}
\affiliation{\institution{AI Coding Agent (GitHub)}\country{}}

\author{Claude Sonnet 4.6}
\affiliation{\institution{AI Model (Anthropic)}\country{}}
\email{}

\begin{abstract}
Formal verification has historically demanded deep expertise in proof
assistants and months of dedicated human effort, limiting its adoption to
a handful of well-funded teams.  We show an example of how LLM-driven automated agentic
workflows can shift this cost nearly entirely to compute.  Using \emph{Lean
Squad}, an open-source workflow that runs as a scheduled GitHub Actions
job, we produce a machine-checked formal specification and verification of core safety
properties of \texttt{raft-rs}---a production Rust implementation of the
Raft distributed consensus protocol---with near-zero human labour.  The
human author required no expertise in Lean, formal methods, or proof
engineering; his role was limited to designing the reusable workflow,
merging pull requests, and writing four brief GitHub issue comments across 129
autonomous agent runs.

The agent autonomously researched the codebase, selected 26 verification
targets, designed a layered proof architecture, wrote all specifications, and completed all proofs:
718~theorems across 79~Lean~4 files (stdlib only, no Mathlib),
culminating in the end-to-end theorem \texttt{raftReachable\_safe}---for
any cluster state reachable via valid Raft transitions from an empty
initial state, no two nodes ever apply different log entries at the same
committed index---with zero \texttt{sorry}.
All five abstract hypotheses of the \texttt{RaftReachable} model are
fully discharged from concrete protocol theorems, leaving no axioms.
One formulation bug was caught by the \texttt{sorry} discipline;
no implementation bugs were found.
An additional 28 correspondence-validation files provide 625
compile-time \texttt{\#guard} assertions confirming that the Lean models
agree with the Rust implementation on concrete inputs.

The total LLM token cost was approximately \$400 (129~runs at
$\sim$\$3 each using Claude Sonnet~4.6 pricing at time of publication).
The Lean Squad workflow is open-source, installs on any GitHub repository
with two commands, and requires only an LLM API token.  The proof
artefacts and full run history are available at
\texttt{dsyme/raft-lean-squad}.
\end{abstract}

\maketitle

\noindent\textbf{Note on authorship and voice.}  Throughout this paper, ``we''
generally reflects the perspective of \emph{LeanSquad}, an autonomous AI
automation workflow implemented using GitHub Agentic Workflows, with
\emph{Copilot CLI} (GitHub) as the underlying AI coding agent using the \emph{Claude Sonnet~4.6} (Anthropic) LLM.  The human
author (Don Syme) took a deliberately \emph{hands-off} approach: he
designed the automated Lean Squad workflow itself (\S\ref{sec:workflow}) and
merged pull requests, but was not involved in target selection,
the layered proof architecture, specification writing, or any actual
proof engineering.  His total intervention over the project's 129~runs
amounted to four brief GitHub issue comments (\S\ref{sec:interventions}).
\medskip


%% ============================================================
\section{Introduction}
%% ============================================================

Distributed consensus algorithms are notoriously subtle.  The Raft protocol
was designed to be understandable~\cite{ongaro2014raft}, yet its correctness
argument spans term management, log matching, leader completeness, and
state-machine safety across multiple interacting invariants.  Production
implementations add further complexity: joint membership changes, in-flight
message tracking, flow control, and optimised quorum computation.

Formal verification of distributed systems is well-established but
typically expensive: IronFleet~\cite{ironfleet} required person-years of
Dafny proof engineering; Woos~et~al.~\cite{vard} verified a simplified
Raft in $\sim$50,000 lines of Coq.  The expertise barrier---fluency in a
proof assistant, invariant discovery, and multi-month proof campaigns---has
limited adoption to a handful of well-funded teams.  Our contribution
addresses this barrier directly, and is different in three ways:

\begin{enumerate}
\item We target a \emph{production} Raft implementation (\texttt{raft-rs},
  used in TiKV~\cite{raftrs}), not a pedagogical model.
\item We use Lean~4 \emph{stdlib only}---no Mathlib---providing a lightweight,
  self-contained verification artefact.
\item The entire effort is \emph{automated}: every target selection,
  specification, proof architecture decision, and proof was produced by an
  LLM-driven FV agent across 129 independent runs, with only four brief
  human interventions (\S\ref{sec:interventions}).  The human author
  required no Lean expertise, no formal-methods background, and invested
  near-zero hours of proof-engineering labour.
\end{enumerate}

The top-level theorem (\S\ref{sec:results}) is:
\begin{lstlisting}[language=Lean4]
theorem raftReachable_safe [DecidableEq E]
    (hd : Nat) (tl : List Nat) (cs : ClusterState E)
    (hvoters : cs.voters = hd :: tl)
    (h : RaftReachable cs) : isClusterSafe cs
\end{lstlisting}

This is machine-checked by Lean~4 with zero \texttt{sorry}.

%% ============================================================
\section{Background}
%% ============================================================

\subsection{The Raft Consensus Protocol}

Raft replicates a state machine through a distributed log.  A cluster of $n$
nodes elects a leader; the leader accepts log entries from clients and
replicates them to followers via \texttt{AppendEntries} RPCs.  An entry is
\emph{committed} once a majority quorum of nodes has acknowledged it.
Safety requires that no two nodes ever apply different entries at the same
log index.

The \texttt{raft-rs} codebase (derived from TiKV's raft-rs~\cite{raftrs})
implements the Consensus Module in Rust: 52 source files covering quorum
arithmetic, log management, flow control, membership change, and configuration
validation.  It supports joint consensus for membership transitions, a
ring-buffer-based in-flight message tracker, and per-node progress tracking.

\subsection{Lean~4}

Lean~4~\cite{lean4} is a dependently typed proof assistant and programming
language.  Proofs are written in a tactic language with automation including
\texttt{omega} (linear arithmetic), \texttt{simp} (rewriting), and
\texttt{decide}/\texttt{native\_decide} (decidable propositions).
We use Lean~4.30.0-rc2, stdlib only.  All proofs are machine-checked by
the Lean kernel; the \texttt{sorry} keyword marks unproved obligations.

%% ============================================================
\section{The Lean Squad Workflow}
\label{sec:workflow}
%% ============================================================

The Lean Squad workflow~\cite{leansquad} is a GitHub Agentic
Workflow---a declarative specification that orchestrates an LLM coding
agent (e.g.\ Copilot CLI or Codex CLI) as a scheduled GitHub Actions job.  It was designed
by the human author and open-sourced as a reusable workflow template; it
is not specific to Raft or to this repository.  Once installed, the
workflow runs autonomously on a schedule (every 8 hours in our
deployment), with no human in the loop for any individual run.

\subsection{Design Principles}

Four design principles underpin the workflow:

\paragraph{Zero human involvement necessary.}
The agent handles the entire FV effort: research, target selection,
proof architecture design, specification writing, and proof engineering.
The human need only merge PRs and occasionally nudge the process.
This shifts the cost of formal verification from scarce human proof-engineering
expertise to LLM compute, dramatically lowering the barrier to entry
and making FV a practical option for any codebase---not just those
with dedicated verification teams.

\paragraph{Optimistic.} The agent does not wait for perfect conditions.
It starts where the code is and finds what can be verified, even if
that means working with approximations or partial proofs.

\paragraph{Incremental.} Each run contributes one small piece.  A spec
with \texttt{sorry} today becomes a proof next week.  A research note
today drives a formal spec next month.  Persistent memory ensures each
run builds on the last.

\paragraph{Findings-first.} A counterexample is a success, not a failure.
Finding a bug through formal verification is exactly the point.

\paragraph{Transparent.} Every PR, issue, and comment includes an AI
disclosure.  Specs document their approximations.
Proofs document what is still open. The history of every run is public, and the
full history of the specifications and proofs is available in the repository for the
human and the agent to examine as needed.

\subsection{Phase-Weighted Task Selection}
\label{sec:tasksel}

Each run begins with a deterministic pre-step that inspects the
repository's FV artefacts (Lean files, spec documents, open issues and
PRs) and computes a phase-weighted probability distribution over
ten tasks.  Two tasks are drawn and communicated to the agent;
the agent confirms them against its persistent memory and executes
them.  A mandatory final task (status-issue update) always runs.

The ten tasks, in order of the typical progression, are:

\begin{enumerate}\setlength\itemsep{0pt}
\item \textbf{Research \& Target Identification.}
  Survey the codebase to identify FV-amenable functions, data structures,
  and algorithms.  Document expected benefit, spec size, proof
  tractability, and recommended approach.  Produces \texttt{RESEARCH.md}
  and \texttt{TARGETS.md}.  If a critique document exists, incorporate
  its feedback to adjust priorities.
\item \textbf{Informal Spec Extraction.}
  For the highest-priority unstarted target, extract a precise informal
  specification by reading the code and inferring the design intention---preconditions, postconditions, invariants, edge cases, concrete examples,
  and inferred intent not explicit in the code.  Flags ambiguities for
  maintainer review.
\item \textbf{Formal Spec Writing (Lean~4).}
  Write Lean~4 type definitions, function stubs, and key \texttt{theorem}
  declarations with \texttt{sorry} as proof bodies.  Must pass
  \texttt{lake build} before opening a PR.
\item \textbf{Implementation Extraction.}
  Translate the relevant Rust implementation into Lean~4 functional
  definitions so propositions can be formally connected to actual code.
  For imperative code, create a pure functional model and document the
  abstraction.
\item \textbf{Proof Assistance.}
  Attempt to prove stated propositions using Lean~4 tactics.  When a
  proof obligation cannot be closed, investigate whether the spec or the
  implementation is wrong.  If a counterexample is found, file a bug
  issue.
\item \textbf{Correspondence Review.}
  Review how faithfully each Lean model corresponds to the Rust source.
  Assess correspondence level (exact, abstraction, approximation, or
  mismatch) and document all divergences.
\item \textbf{Proof Utility Critique.}
  Honestly assess whether the FV work is useful---are the proved
  properties meaningful, at the right level of abstraction, and likely
  to catch real bugs?  Identify highest-value gaps.
\item \textbf{CI Automation.}
  Set up and maintain \texttt{lean-ci.yml} to automatically verify
  Lean proofs on every PR and push.
\item \textbf{Project Report.}
  Create and maintain \texttt{REPORT.md}---a comprehensive report with
  mermaid diagrams summarising the entire FV effort.
\item \textbf{Conference Paper.}
  Write and maintain a LaTeX paper under
  \texttt{formal-verification/paper/}, compile to PDF, and commit to
  the repository.
\end{enumerate}

The weighting scheme adapts automatically to the current state of
the repository.  When no FV work exists, Task~1 (Research) dominates.
Once research is done, Task~2 rises.  As informal specs accumulate,
Task~3 gains weight, and so on up through proofs and beyond to
reflective tasks (critique, correspondence review, paper writing).
Table~\ref{tab:weights} illustrates the weight progression.

\begin{table}[h]
\centering\small
\begin{tabular}{clcc}
\toprule
Task & Name & Weight (early) & Weight (late) \\
\midrule
1 & Research         & 10.0 & 2.0 \\
2 & Informal Spec    &  8.0 & 2.0 \\
3 & Formal Spec      &  8.0 & 2.0 \\
4 & Implementation   &  6.0 & 3.0 \\
5 & Proof Assistance &  6.0 & 5.0 \\
6 & Correspondence   & 10.0 & 5.0 \\
7 & Critique         & 10.0 & 5.0 \\
8 & CI               & 12.0 & 1.0 \\
9 & Report          &  8.0 & 3.0 \\
10 & Paper           &  8.0 & 3.0 \\
\bottomrule
\end{tabular}
\caption{Phase-weighted task selection.  ``Early'' weights apply when the
  task's prerequisites are met but the task has not yet been completed;
  ``late'' weights apply once the artefact exists and is being maintained.}
\label{tab:weights}
\end{table}

\subsection{Persistent Memory and Cumulative Progress}

A critical design element is persistent repo-memory.  At the start of
every run, the agent reads its memory to recover the identified FV
targets, their current phases (1--5), key choices, open questions, bugs
found, and ideas to try.  At the end of every run, the agent updates
memory with new discoveries, proved theorems, and remaining obligations.

Memory may be stale (referencing merged PRs or closed issues), so the
agent is instructed to verify references against the current repository
state.  This mechanism enables cumulative progress across independent
runs: the layered architecture, the choice of targets, and the
proof strategy all emerged organically from the agent's interaction
with its own memory and the growing codebase, not from human direction.

\subsection{Hard Requirements}

The workflow enforces several hard requirements to maintain integrity:

\begin{itemize}\setlength\itemsep{0pt}
\item The Lean toolchain \textbf{must} be successfully installed before
  any \texttt{.lean} files are written or submitted.  If installation
  fails, Tasks~3--5 are skipped entirely for that run.
\item Every PR that includes \texttt{.lean} files must include a
  verification status block: either ``\texttt{lake build} passed'' or
  an explicit disclosure that proofs are unverified.
\item The agent must never claim proofs are verified unless
  \texttt{lake build} actually passed.
\item All PRs, issues, and comments must include an AI transparency
  disclosure.
\end{itemize}

\subsection{Artefacts Produced}

The workflow produces a rich set of artefacts beyond Lean proofs:
\texttt{RESEARCH.md} (FV survey and rationale),
\texttt{TARGETS.md} (prioritised target list with phase status),
per-target informal specs,
\texttt{CORRESPONDENCE.md} (Lean-to-Rust mapping),
\texttt{CRITIQUE.md} (proof utility assessment),
\texttt{REPORT.md} (comprehensive project report),
and the conference paper itself.  A continuously-updated GitHub
issue serves as a rolling dashboard showing the status of every
target.

\subsection{Installation and Usage}
\label{sec:install}

The Lean Squad workflow is open-source and can be installed on any
GitHub repository with two commands.  A token for an LLM coding agent
is required; at the time of writing, Copilot CLI (GitHub),
Codex CLI (OpenAI), and Claude Code (Anthropic) are supported as
underlying agents.

\begin{lstlisting}[language={},basicstyle=\ttfamily\scriptsize,
  keywordstyle={},frame=single,xleftmargin=4pt,xrightmargin=4pt]
# Install the GitHub Agentic Workflows CLI extension
gh extension install github/gh-aw

# Add Lean Squad to the current repository
gh aw add-wizard githubnext/agentics/lean-squad
\end{lstlisting}

The \texttt{add-wizard} command walks the user through configuration:
selecting the schedule (default: every 8 hours), choosing the agent
provider, and storing the API token as a GitHub Actions secret.
Once committed, the workflow runs autonomously on the configured
schedule with no further setup.

With just these two commands, any repository can start an automated formal verification
effort with zero human specification or proof-engineering expertise required.
Whether that is useful or not depends on the codebase and the agent's
performance, but the barrier to entry is minimal. Getting started with
formal verification of real software has never been easier.

%% ============================================================
\section{Methodology}
%% ============================================================

\subsection{Target Selection}
\label{sec:targets}

In its very first run (PR~\#93), the agent surveyed the
\texttt{raft-rs} codebase via Task~1 (Research) and identified an
initial set of \textbf{10~FV targets}, prioritised by tractability:
\texttt{limit\_size}, \texttt{Config::validate}, \texttt{vote\_result},
\texttt{committed\_index}, \texttt{find\_conflict},
\texttt{maybe\_append}, joint-config variants, \texttt{inflights}, and
\texttt{progress}.  That same run also produced the first informal
specification (\texttt{limit\_size}).  The criteria were: pure
functions with clear pre/postconditions, data-structure invariants,
and proof tractability.  No human input was involved in selecting or
prioritising targets.

Notably, the first run's research plan proposed Lean~4 + Mathlib and
flagged the Aeneas extraction toolchain as a possible route; both
choices were later revised by the agent itself---switching to
stdlib-only and dropping Aeneas---as it gained experience with the
proof artefacts across subsequent runs.

Over later iterations the agent expanded the target list to
\textbf{26~targets} and organised them into seven layers
(Table~\ref{tab:layers}).  Twelve additional targets were added as
the agent discovered cross-module composition opportunities
(\texttt{SafetyComposition}, \texttt{CrossModuleComposition}),
protocol-level properties (\texttt{RaftSafety}, \texttt{RaftProtocol}),
the reachability model (\texttt{RaftTrace}), further quorum
lemmas (\texttt{HasQuorum}, \texttt{TallyVotes},
\texttt{QuorumRecentlyActive}, \texttt{IsUpToDate},
\texttt{LogUnstable}), and the election model
(\texttt{RaftElection}, \texttt{LeaderCompleteness},
\texttt{ConcreteTransitions}, \texttt{CommitRule},
\texttt{MaybeCommit}, \texttt{ConcreteProtocolStep}).
This layered architecture emerged from the
agent's own analysis of module dependencies---it was not prescribed by
the human author.
Layers build on each other: Layer~2 (Quorum Arithmetic) uses types from
Layer~1; Layer~4 (Composition) spans Layer~2 and Layer~3; Layers~5--6
depend on all preceding layers; Layer~7 (Election Model) provides
concrete discharge lemmas for the abstract step hypotheses of Layer~6.

\begin{table}[h]
\centering\small
\begin{tabular}{clrr}
\toprule
Layer & Description & Files & Theorems \\
\midrule
\multicolumn{4}{l}{\emph{Core proof architecture (Layers 1--7):}} \\
1 & Data Structures     & 6  & 161 \\
2 & Quorum Arithmetic   & 8  & 135 \\
3 & Protocol Operations & 3  &  43 \\
4 & Cross-Module Composition & 3 & 26 \\
5 & Raft Safety         & 2  & 18  \\
6 & Reachability        & 2  & 11  \\
7 & Election Model (complete) & 12 & 152 \\
\midrule
\multicolumn{4}{l}{\emph{Correspondence validation (Layer 8):}} \\
8 & Correspondence Validation & 28 & 2+625\texttt{g} \\
\midrule
\multicolumn{4}{l}{\emph{Additional module verifications (Layers 9--20):}} \\
9  & ReadOnly Operations      & 1  & 14  \\
10 & FindConflictByTerm       & 1  & 10  \\
11 & Flow Control             & 1  & 18  \\
12 & Configuration Invariants & 1  & 11 \\
13 & Persistent-Storage Bridge & 2 & 16 \\
14 & Progress Tracking        & 1  & 26 \\
15 & Raft-Log Append          & 1  & 16 \\
16 & Next-Entries Slicing     & 2  & 21 \\
17 & Progress Set             & 1  &  8 \\
18 & MemStorage Core          & 1  & 14 \\
19 & IsContinuousEnts         & 1  &  8 \\
20 & Restore                  & 1  &  8 \\
\midrule
  & \textbf{Total}      & \textbf{79} & \textbf{718+625\texttt{g}} \\
\bottomrule
\end{tabular}
\caption{Proof layers and theorem counts.  Layers 1--7 form the
core safety-proof chain; Layers 9--20 are standalone module
verifications.  Layer~8 files contain 625~\texttt{\#guard} assertions
(\texttt{g}) evaluated at compile time; they contribute 2 named
theorems (in \texttt{FindConflictCorrespondence.lean}).
File count includes \texttt{Basic.lean} (helper definitions, 0
theorems).}
\label{tab:layers}
\end{table}

\subsection{Specification Strategy}

Each target was formalised in two phases.  First, an \emph{informal
specification} was extracted by reading the Rust source, its tests, and
documentation comments.  This produced a prose document identifying
preconditions, postconditions, invariants, edge cases, and open questions.

Second, the informal spec drove a \emph{Lean~4 formalisation}: type
definitions mirroring the Rust types (with deliberate approximations
documented in CORRESPONDENCE.md), function stubs, and theorem statements.
Proofs were then attempted incrementally, from simplest (structural
induction, \texttt{omega}) to most complex (multi-lemma composition).

\subsection{Modelling Choices}
\label{sec:modelling}

The Lean model deliberately abstracts several Rust features:

\paragraph{Type abstraction.} \texttt{u64} voter IDs and index values
are modelled as \texttt{Nat}.  This eliminates overflow concerns (safe: no
practical Raft deployment approaches $2^{64}$) and enables \texttt{omega}
arithmetic.

\paragraph{Collection abstraction.} \texttt{HashSet<u64>} voter sets are
modelled as \texttt{List Nat}.  Theorems that require uniqueness carry an
explicit \texttt{voters.Nodup} hypothesis; those that do not hold for
arbitrary lists.

\paragraph{Log abstraction.} The per-node log is modelled as
\texttt{VoterLogs E := Nat $\to$ Nat $\to$ Option E}---a pure function
from (voter, index) to optional entry.  This abstracts over the Rust
\texttt{Vec}-based log structure while preserving the read semantics.

\paragraph{Protocol abstraction.} Concrete Raft message types are
modelled abstractly: \texttt{ValidAEStep} captures AppendEntries effects,
\texttt{ElectionEpoch} captures election outcomes, and
\texttt{BroadcastSeq} captures broadcast rounds.  Term management and
leader election are formalised at the level needed to discharge the
\texttt{RaftReachable.step} hypotheses (Layer~7), but the full
message-passing network model (UDP delivery, reordering, partitions)
is not included.

\paragraph{Size semantics.} For \texttt{limit\_size}, protobuf message
size is abstracted as an arbitrary \texttt{size : $\alpha \to$ Nat} function.
Proofs hold for any such function, making them strictly more general than
the Rust implementation.

\subsection{Proof Approach}

We used the following tactic strategy, from simplest to most complex:

\begin{itemize}\setlength\itemsep{0pt}
  \item \texttt{omega}: linear integer/natural arithmetic (most common)
  \item \texttt{simp}/\texttt{simp only}: definitional unfolding and rewriting
  \item \texttt{by\_cases}: boolean and decidable case splits
  \item \texttt{induction}/\texttt{cases}: structural recursion on lists and inductives
  \item \texttt{exact}/\texttt{apply}/\texttt{refine}: explicit term construction
  \item \texttt{funext}: function extensionality
\end{itemize}

We used no \texttt{native\_decide}, no axioms beyond Lean's kernel, and
no \texttt{sorry} in the final proof state.

\subsection{Model Validation (Layer 8)}

A proof that \emph{the Lean model} is correct conveys confidence only if
the Lean model faithfully captures the Rust implementation.
To ground the abstract models, we developed 26 \emph{correspondence-test}
files (Layer~8) containing compile-time \texttt{\#guard} assertions.  Each
file re-evaluates the Lean model function on concrete inputs and checks
the result against the expected value from the Rust implementation.
Because \texttt{\#guard} is evaluated by \texttt{lake build} at compile time
(not at runtime), these are not optional tests---they are compile-time
obligations that must pass for the build to succeed.

For each target, a matching Rust function (prefixed \texttt{test\_\ldots\_correspondence})
exercises the Rust implementation on the same concrete inputs, providing
independent confirmation.  Together, the 625 Lean \texttt{\#guard}
assertions and 18+ Rust test functions provide two-sided behavioural
evidence: both the Lean model and the Rust implementation produce
identical outputs on the verified test cases.

%% ============================================================
\section{Results}
%% ============================================================

\label{sec:results}

\subsection{Layer-by-Layer Proof Inventory}

\paragraph{Layer~1: Data Structures.}
Six files prove correctness of individual Raft data structures in isolation.

\texttt{LimitSize.lean} (17~theorems) formalises the Rust function
\texttt{util::limit\_size}, which truncates a vector of protobuf
messages to stay within a byte budget.
Key results: \emph{prefix correctness} (the result is a prefix of
the input), \emph{non-emptiness} ($\geq 1$ entry always preserved), and
\emph{maximality} (the longest prefix within budget).

\texttt{Inflights.lean} (40 theorems) is the largest single file.  It
formalises the ring-buffer-based in-flight message tracker (\texttt{inflights.rs}),
proving the abstract/concrete bridge: the efficient ring buffer and a
simple list model compute identical \texttt{free\_to} results.

\texttt{Progress.lean} (36 theorems) formalises the per-node progress
tracker state machine, proving the well-foundedness invariant
\texttt{matched + 1 $\leq$ next\_idx} is preserved across all transitions
(including \texttt{maybeDecrTo} probe step-back and \texttt{optimisticUpdate} pipeline advance).

\paragraph{Layer~2: Quorum Arithmetic.}
Eight files prove the mathematical foundations of Raft safety.
The cornerstone is \texttt{quorum\_intersection\_mem} (HQ20,
\texttt{HasQuorum.lean}):

\begin{lstlisting}[language=Lean4]
theorem quorum_intersection_mem
  (voters : List Nat) (hne : voters ≠ [])
  (a b : Nat → Bool)
  (ha : hasQuorum voters a = true)
  (hb : hasQuorum voters b = true) :
  ∃ v, v ∈ voters ∧ a v = true ∧ b v = true
\end{lstlisting}

This states: any two majority quorums over the same voter list share at
least one concrete member.  This is the quorum-intersection property that
makes Raft (and Paxos) safe.  The proof proceeds by an arithmetic argument
using inclusion-exclusion and \texttt{omega}.

\texttt{CommittedIndex.lean} (17 theorems) formalises the
\emph{sort-then-index} algorithm that computes the committed index from
a list of per-voter match indices, proving both safety (result is
attested by a majority) and maximality (result is the largest such index).

\paragraph{Layer~3: Protocol Operations.}
Three files prove correctness of core Raft log operations.
\texttt{MaybeAppend.lean} (18 theorems) gives a complete post-condition
for \texttt{maybe\_append}: theorem MA13 proves the prefix is untouched
and MA14 proves the suffix is correctly applied.  Together they provide
a behavioural specification of AppendEntries log modification.

\texttt{FindConflict.lean} (12 theorems) proves that the function
returns zero iff all entries match (FC11), and returns the first index
of mismatch otherwise.

\texttt{MaybePersist.lean} (13 theorems) formalises the async-persist
safety guards of \texttt{maybe\_persist} and \texttt{maybe\_persist\_snap}
from \texttt{src/raft\_log.rs}.  Key results: MP1 gives the complete
iff-characterisation of the three guards (monotonicity, below-firstUpdateIndex,
term-match); MP8 proves idempotence (second call with the same arguments
always returns false, preventing double-advance).

\paragraph{Layer~4: Cross-Module Composition.}
Three files prove properties spanning multiple independent modules.
\texttt{SafetyComposition.lean} (9 theorems) bridges quorum arithmetic
and committed index computation: theorem SC4 states that if two voters
have committed different entries at the same index, they must share a
witness voter in their respective quorums---but the same voter cannot
have voted for two different entries, giving a contradiction.

\texttt{CrossModuleComposition.lean} (7 theorems) proves that
\texttt{maybe\_append} is safe with respect to quorum commitments (CMC3):
it never advances a committed index beyond what the current quorum has
certified.

\paragraph{Layer~5: Raft Safety.}
\texttt{RaftSafety.lean} (10 theorems) lifts quorum-level safety to
the log-entry level.  The central theorem is RSS1:

\begin{lstlisting}[language=Lean4]
theorem raft_state_machine_safety
  [DecidableEq E]
  (voters : List Nat) (hne : voters ≠ [])
  (logs : VoterLogs E) (k : Nat)
  (e₁ e₂ : E)
  (h1 : isQuorumCommitted voters logs k e₁)
  (h2 : isQuorumCommitted voters logs k e₂)
  : e₁ = e₂
\end{lstlisting}

If two entries are simultaneously quorum-committed at the same index,
they must be identical.  The proof uses \texttt{quorum\_intersection\_mem}
(HQ20) to obtain a shared witness voter and conclude that both entries
equal the witness's log entry.

RSS8 (\texttt{raft\_end\_to\_end\allowbreak\_safety\_full}) provides the
conditional end-to-end result: given the
\texttt{CommitCertInvariant}---every applied
entry is quorum-certified---the cluster is safe.

\texttt{RaftProtocol.lean} (8~theorems) proves that
\texttt{Append\-Entries} preserves the Log Matching Invariant (RP6) and
the No-Rollback property (RP8), formalising the Raft protocol transition
rules at the log level.

\paragraph{Layer~6: Reachability.}

\begin{figure}[h]
\begin{lstlisting}[language=Lean4]
inductive RaftReachable [DecidableEq E] :
  ClusterState E → Prop where
| init (voters : List Nat) :
    RaftReachable (initialCluster voters)
| step {cs cs' : ClusterState E}
    (hreach : RaftReachable cs)
    (hvoters : cs.voters = cs'.voters)
    (v : Nat)
    (hlogs' : ∀ w, w ≠ v →
      cs'.logs w = cs.logs w)
    (hno_overwrite : ∀ k,
      cs.committed v ≥ k →
      cs'.logs v k = cs.logs v k)
    (hqc_preserved : ∀ k e,
      isQuorumCommitted
        cs.voters cs.logs k e →
      ∀ w, cs'.logs w k = cs.logs w k)
    (hcommitted_mono : ∀ w,
      cs'.committed w ≥ cs.committed w)
    (hnew_cert : ∀ w k e,
      cs'.committed w ≥ k →
      cs.committed w < k →
      cs'.logs w k = some e →
      isQuorumCommitted
        cs'.voters cs'.logs k e) :
    RaftReachable cs'
\end{lstlisting}
\caption{The \texttt{RaftReachable} inductive protocol model.}
\label{fig:rr}
\end{figure}

\texttt{RaftTrace.lean} (4~theorems) closes the gap between the conditional
safety result (RSS8) and an unconditional statement.  It defines the
\texttt{RaftReachable} inductive predicate (Figure~\ref{fig:rr}) and
proves by structural induction that every reachable state satisfies
\texttt{CommitCertInvariant} (RT1), from which RSS8 yields the top-level
theorem RT2:

\begin{lstlisting}[language=Lean4]
theorem raftReachable_safe [DecidableEq E]
    (hd : Nat) (tl : List Nat) (cs : ClusterState E)
    (hvoters : cs.voters = hd :: tl)
    (h : RaftReachable cs) : isClusterSafe cs :=
  raft_end_to_end_safety_full hd tl cs hvoters
    (raftReachable_cci cs h)
\end{lstlisting}

The two-line proof follows immediately from RT1 and RSS8.

\texttt{MultiStepReachability.lean} (7~theorems) extends this to $n$-step
chains: theorem MS2 provides an end-to-end safety certificate showing that
a valid $n$-step AE sequence from a \texttt{RaftReachable} state yields
another \texttt{RaftReachable} state satisfying the safety invariant.

\paragraph{Layer~7: Election Model (complete).}

Twelve additional files formalise Raft's election and commit machinery
at a more concrete level, providing the theorems that fully discharge the
\texttt{RaftReachable.step} hypotheses.

\texttt{RaftElection.lean} (53~theorems) formalises abstract Raft-step
properties including single-voter log-change invariants, providing
lemmas used by higher layers.

\texttt{LeaderCompleteness.lean} (10~theorems) formalises leader
completeness properties needed to discharge \texttt{hqc\_preserved}: the
property that a newly elected leader has all quorum-certified entries
in its log.

\texttt{ConcreteTransitions.lean} (8~theorems) formalises the
\texttt{AppendEntries} effect function directly, proving CT1--CT5b:
key properties of how the log is updated by a concrete AE step.

\texttt{CommitRule.lean} (9~theorems) formalises the Raft commit rule
(\texttt{leaderCommit} reflects a genuine quorum acknowledgement).
Theorem CR8 proves \texttt{CommitRuleValid~=~hnew\_cert} by
definitional equality (\texttt{Iff.rfl}), \emph{closing} the
\texttt{hnew\_cert} hypothesis of \texttt{RaftReachable.step}.

\texttt{MaybeCommit.lean} (12~theorems) formalises \texttt{maybeCommit},
the function that advances the commit index.  Theorem MC4
(\texttt{maybeCommit\_term}) proves \emph{A6 term safety}: the commit index
advances only when the entry at the new committed index has
\texttt{term~=~cs.term} (leader's current term), formalising Raft~\S5.4.2.
This is the formal Lean proof of the mechanism that prevents the ``figure~8''
problem from the Raft paper: a stale entry from an old leader's term cannot
be directly committed by a new leader, because \texttt{maybeCommit}
would use the new leader's current term, and the old entry would carry a
different term.  MC4 and CR8 together close both halves of commit
safety: CR8 proves the quorum-certification half, MC4 proves the term-safety half.

\texttt{ConcreteProtocolStep.lean} (14~theorems) introduces the
\texttt{ValidAEStep} structure and proves the \emph{A5 bridge}
(CPS2, \texttt{validAEStep\_raftReachable}): if \texttt{cs} is
\texttt{RaftReachable} and a valid \texttt{ValidAEStep} fires on it,
the resulting state \texttt{cs'} is also \texttt{RaftReachable}.  This is
the first theorem that directly connects a concrete Raft message to the
abstract reachability model.

Four additional files in Layer~7 close the \texttt{hqc\_preserved} gap:

\texttt{ElectionReachability.lean} (12~theorems) proves that
\texttt{CandidateLogCovers} holds under the \emph{high-water-mark}
condition: if a reference log exists from which all voters received their
entries, the candidate's log covers every committed entry.  Theorems
ER1--ER4 reduce the obligation; ER9--ER10 close it under the
shared-source condition.

\texttt{ElectionConcreteModel.lean} (8~theorems) bridges the abstract
hypotheses of \texttt{ElectionReachability} to the concrete Raft protocol:
ECM1--ECM4 prove that a \texttt{ValidAEStep} establishes the
\texttt{hae} (log-agreement after AE) hypothesis, and ECM5--ECM7 chain
this through \texttt{CandidateLogCovers} to discharge
\texttt{hqc\_preserved}.

\texttt{AEBroadcastInvariant.lean} (10~theorems) proves that the
log-agreement invariant \texttt{hae} holds inductively as the leader
broadcasts AppendEntries to voters one by one.  The key theorem,
ABI8 (\texttt{hqc\_preserved\_of\_broadcast}), establishes that after a
complete AE broadcast round, the cluster satisfies \texttt{hqc\_preserved}.

\texttt{ElectionBroadcastChain.lean} (6~theorems) closes the gap left by
the \texttt{haeCovered\_induction} placeholder in
\texttt{AEBroadcastInvariant.lean}.  It introduces \texttt{BroadcastSeq},
an inductive type threading cluster states through a sequence of
single-voter AppendEntries steps.  EBC1--EBC2 prove log-stability for
non-stepping nodes; EBC3 proves per-voter agreement through the tail;
EBC4 (\texttt{broadcastSeq\_haeCovered}) is the key structural-induction
theorem establishing \texttt{HaeCovered} for the entire voter list;
EBC5 lifts \texttt{HaeCovered} to a universal \texttt{hae} hypothesis;
and EBC6 (\texttt{electionBroadcastChain\_hqc\_preserved}) chains these
through ECM4 (\texttt{hqc\_preserved\_of\_hae}) to deliver
\texttt{hqc\_preserved} from a complete broadcast sequence---the concrete
bridge previously missing between the AE broadcast invariant and
\texttt{hqc\_preserved}.

\texttt{ElectionLifecycle.lean} (7~theorems, EL1--EL7) is the capstone
file that closes the gap.  It defines \texttt{ElectionEpoch}---a
structure tying an election winner to a specific term and broadcast
round---and proves that any \texttt{RaftReachable} state followed by an
\texttt{ElectionEpoch} and a \texttt{ValidAEStep} produces a new
\texttt{RaftReachable} state.  EL1 (\texttt{electionEpoch\_hqc\_pre\-served})
discharges the hardest hypothesis; EL7 (\texttt{fullProtocol\-Step\_safe})
is the end-to-end result: no abstract axioms remain.

\texttt{RaftLog\-Append.lean} (16~theorems) formalises
\texttt{raft\_log.rs::append} as a Lean~4 function over a
\texttt{RaftLog} record, proving
RA1--RA9: the empty-input invariant, the return-value specification
(last index after append), the committed-index monotonicity, and the
stable-storage prefix preservation.  Theorems P4--P7 prove further
prefix and prefix-of-committed properties.

\texttt{MaybePersistFUI.lean} (8~theorems) formalises the
\texttt{firstUpdateIndex} derivation in \texttt{RaftLog::maybe\_persist}
(\texttt{src/raft\_log.rs}).  The function determines the boundary below
which persistent storage must not advance, derived from the \texttt{Unstable}
struct's snapshot and entries fields.  Key theorems: FU1
(\texttt{fu1\_unstable\_snap\_blocks}): when a snapshot is set, the FUI
equals \texttt{snap.index + 1}, fencing off persistence below the snapshot
horizon; FU8 (\texttt{fu8\_unstable\_fui\_gt\_persisted}): the derived FUI
is strictly greater than the current \texttt{persisted} value---the key
safety invariant ensuring \texttt{maybe\_persist} cannot advance past stable
storage.  All 8 theorems have 0 \texttt{sorry}.

\paragraph{Layer~8: Correspondence Validation.}
Twenty-eight correspondence-test files provide compile-time evidence that the
Lean models agree with the Rust implementation on concrete inputs.  Each
file re-implements a specific Lean model function and checks it against
hand-crafted fixtures via \texttt{\#guard} assertions, which are evaluated
by \texttt{lake build} at compile time.  Matching Rust test functions run
the same cases via \texttt{cargo test}.

The 28 files cover: \texttt{find\_conflict} (20~guards, exact),
\texttt{maybe\_append} (32~guards, exact), \texttt{is\_up\_to\_date}
(12~guards, exact), \texttt{committed\_index} (11~guards, abstraction),
\texttt{limit\_size} (10~guards, abstraction), \texttt{config\_validate}
(12~guards, exact), \texttt{inflights} (12~guards, abstraction),
\texttt{log\_unstable} (12~guards, exact), \texttt{tally\_votes}
(10~guards, exact), \texttt{vote\_result} (15~guards, exact),
\texttt{has\_quorum} (15~guards, exact),
\texttt{read\_only} (15~guards, exact),
\texttt{find\_conflict\_by\_term} (16~guards, abstraction),
\texttt{progress} (51~guards, abstraction --- \texttt{maybeUpdate},
\texttt{maybeDecrTo}, \texttt{optimisticUpdate}, and state transitions),
\texttt{progress\_tracker} (44~guards, abstraction),
\texttt{maybe\_persist} (21~guards, abstraction --- \texttt{maybe\_persist}
and \texttt{maybe\_persist\_snap} with snapshot fixtures),
\texttt{maybe\_commit} (17~guards, abstraction --- \texttt{maybe\_commit}
and \texttt{commit\_to}: the A6 term-safety guard and monotone commit),
\texttt{raft\_log\_append} (21~guards, abstraction --- all three
\texttt{truncate\_and\_append} branches),
\texttt{maybe\_persist\_fui} (25~guards, abstraction --- FUI derivation),
\texttt{election} (28~guards, abstraction),
\texttt{configuration\_invariants} (18~guards, abstraction),
\texttt{joint\_vote} (18~guards, abstraction),
\texttt{quorum\_recently\_active} (14~guards, abstraction),
\texttt{uncommitted\_state} (13~guards, abstraction),
\texttt{has\_next\_entries} (33~guards, abstraction),
\texttt{next\_entries} (28~guards, abstraction),
\texttt{progress\_set} (28~guards, abstraction), and
\texttt{mem\_storage} (23~guards, abstraction).
In total: 625 compile-time correspondence assertions and 18+ Rust test functions.

\paragraph{Layer~9: ReadOnly Operations.}
\texttt{ReadOnly.lean} (14~theorems) formalises \texttt{raft-rs}'s
read-only request queue (\texttt{src/read\_only.rs}), modelling it as a
pure record with a request queue and a pending-ack map.  Key theorems include
RO8 (\texttt{advance\_removes\_ctx}: advancing the committed index dequeues
exactly the confirmed read requests) and RO14--RO15 (the \texttt{QueuePendingInv}
invariant is preserved by \texttt{recvAck} and \texttt{advance}).

\paragraph{Layer~10: Backward-Scan Optimisation (FindConflictByTerm).}
\texttt{FindConflictByTerm.lean} (10~theorems) formalises
\texttt{find\_conflict\_by\_term} (\texttt{src/raft\_log.rs},
lines 218--257),
which allows a Raft leader to skip over a diverging log suffix in a
single step given a conflict hint from a rejecting follower.
Key theorems: FCB1 (result index is bounded by input), FCB3 (maximality
of the backward scan --- every skipped position has term exceeding the
hint), FCB6 (the scan always terminates with a \texttt{some} result via
the dummy entry at index~0), and FCB9 (with a non-decreasing log the
scan finds the largest valid index).

\paragraph{Layer~11: Flow-Control Byte Accounting (UncommittedState).}
\texttt{UncommittedState.lean} (18~theorems) formalises the
\texttt{UncommittedState} struct (\texttt{src/raft.rs}) responsible for
flow-controlling uncommitted-log byte usage.  Key theorems: US16
(complete characterisation of budget rejection), US17 (dual: acceptance
iff at least one free-pass condition holds), and US18 (saturation-safe
lower bound on \texttt{maybeReduce} output).

\paragraph{Layer~12: Configuration Invariants.}
\texttt{Configuration\-Invariants.lean} (11~theorems) formalises the
voter/learner disjointness invariant across joint-consensus
configuration changes.  Key theorems: CI5 (\texttt{allVoters} $\cap$ learners = $\emptyset$
under the invariant), CI7 (invariant preserved under joint-config
finalisation), and CI11 (counterexample showing the strict invariant is
strictly stronger than the relaxed version during demotion state).

\paragraph{Layer~13: Persistent-Storage Bridge.}
\texttt{MaybePersistFUI.lean} (8~theorems) and
\texttt{UnstablePersistBridge.lean} (8~theorems) together bridge the
\texttt{maybe\_persist} operation to the \texttt{LogUnstable}
well-formedness predicate.  FU1 shows that a pending snapshot blocks
the firstUpdateIndex from advancing; FU8 proves fui $>$ persisted as a
safety invariant; UPB8 closes the async-persist safety gap.

\paragraph{Layer~14: Progress Tracking.}
\texttt{ProgressTracker.lean} (26~theorems, PT1--PT26) proves that
the \texttt{all\_wf} invariant---every peer's progress record is
well-formed---is preserved by all six progress-tracker operations:
\texttt{recordVote}, \texttt{quorumRecentlyActive}, \texttt{isSingleton},
\texttt{applyChanges}, and membership update, plus two integration
theorems (PT25--PT26) connecting \texttt{initTracker} and
\texttt{applyChanges} to membership preservation.

\paragraph{Layer~15: Raft-Log Append.}
\texttt{RaftLog\-Append.lean} (16~theorems, RA1--RA9 + P4--P7) formalises
\texttt{RaftLog::\allowbreak truncate\_and\_append}
(\texttt{src/raft\_log.rs}).
Key results: RA4 (prefix preservation under truncation), RA7 (suffix
correctly applied), and P4--P7 (stable-storage prefix preservation
across append).

\paragraph{Layer~16: Next-Entries Slicing.}
\texttt{NextEntries.lean} (7~theorems, NE1--NE7) formalises
\texttt{RaftLog::next\_entries\_since} and \texttt{next\_entries}
(\texttt{src/raft\_log.rs\#L447--L492}).  Key results: NE1 (returns
\texttt{none} iff \texttt{hasNextEntriesSince} is false), NE3 (the
$i$-th element equals $\mathit{offset}+i$), NE4 (anti-monotone in
\texttt{sinceIdx}), and NE7 (\texttt{next\_entries} equals
\texttt{next\_entries\_since(applied)}).
\texttt{HasNextEntries.lean} (14~theorems, HNE1--HNE14) formalises the
boolean predicate \texttt{hasNextEntries\-Since} / \texttt{hasNext\-Entries}.
HNE7 characterises when the predicate is false (empty slice); HNE14 gives the
top-level conjunction form used by NE1.

\paragraph{Layer~17: Progress Set.}
\texttt{ProgressSet.lean} (8~theorems, PS1--PS8) formalises the
\texttt{ProgressTracker} as a composed progress set and proves 8 key
properties of \texttt{quorum\_recently\_active} as a composed operation.
PS1 proves the quorum threshold; PS8 proves the false-negative freedom.
Correspondence validated by \texttt{ProgressSetCorrespondence.lean} (28
\texttt{\#guard}).

\paragraph{Layer~18: MemStorage Core.}
\texttt{MemStorage.lean} (14~theorems) models \texttt{MemStorageCore}
(\texttt{src/storage.rs}) pure log operations: \texttt{firstIndex},
\texttt{lastIndex}, \texttt{compact}, and \texttt{append}.  Key results:
contiguity preservation under compact, append length correctness, and
monotonicity of \texttt{firstIndex} under compaction.  Correspondence
validated by \texttt{MemStorageCorrespondence.lean} (23~\texttt{\#guard}).

\paragraph{Layer~19: IsContinuousEnts.}
\texttt{IsContinuousEnts.lean} (8~theorems, ICE1--ICE8) formalises
\texttt{is\_continuous\_ents} (\texttt{src/util.rs}), which validates that two
AppendEntries batches can be safely concatenated.  ICE1--ICE3 cover the
three branches (empty msg, empty ents, index check); ICE4--ICE8 prove
structural properties including symmetry and transitivity of the predicate.

\paragraph{Layer~20: Restore.}
\texttt{Restore.lean} (8~theorems, RST1--RST8) formalises
\texttt{RaftLog::restore} (\texttt{src/raft\_log.rs:696})---the snapshot
restoration operation.  RST1 (committed = snapshot index), RST2 (persisted
$\le$ committed), RST3 (committed advances), RST7 (ordering invariant
preserved), and RST8 (idempotence at current committed) are the key results.

\subsection{Proof Statistics}

Table~\ref{tab:inventory} summarises the full file inventory.
Total: 718 theorems, 79~files (50 proof files + 28 correspondence-test
files + 1 helper), 0 \texttt{sorry}, $\approx$21{,}300 lines of Lean.

\begin{table*}[p]
\centering\footnotesize
\begin{tabular}{lrp{5cm}}
\toprule
File & Thms & Key result \\
\midrule
LimitSize         & 17 & Prefix, maximality \\
ConfigValidate    & 10 & 8-constraint predicate equiv \\
MajorityVote      & 21 & Won/Lost/Pending characterisation \\
JointVote         & 14 & Joint-config degeneracy \\
CommittedIndex    & 17 & Sort-index safety + maximality \\
FindConflict      & 12 & First mismatch; zero $\Leftrightarrow$ all-match \\
JointCommittedIndex & 10 & $\min(\mathit{ci}_{in}, \mathit{ci}_{out})$ \\
MaybeAppend       & 18 & Prefix preserved; suffix applied \\
Inflights         & 40 & Ring-buffer abstract/concrete bridge \\
Progress          & 36 & State-machine WF invariant \\
IsUpToDate        & 17 & Lex order, total preorder \\
LogUnstable       & 41 & Snapshot/entries consistency \\
TallyVotes        & 28 & Partition, bounds, no-double-quorum \\
HasQuorum         & 20 & Quorum intersection (HQ14, HQ20) \\
QuorumRecentlyActive & 11 & Active-quorum detection \\
SafetyComposition & 9  & SC4: two CIs share witness \\
JointTally        & 14 & Joint-config tally correctness \\
JointSafetyComposition & 10 & JSC7: witnesses in both halves \\
CrossModuleComposition &  7 & CMC3: maybe\_append bounded \\
RaftSafety        & 10 & RSS1--RSS8: cluster safety \\
RaftProtocol      &  8 & LMI, NRI preserved by AppendEntries \\
RaftTrace         &  4 & RT1, RT2: reachable safety \\
RaftElection      & 53 & Abstract Raft-step properties (incl.\ election) \\
LeaderCompleteness & 10 & Leader completeness components \\
ConcreteTransitions &  8 & CT4: LMI preserved; CT5: broadcast $\to$ match \\
CommitRule        &  9 & CR8: \texttt{hnew\_cert} closed \\
MaybeCommit       & 12 & MC4: A6 term safety \\
MaybePersist      & 13 & MP6: async-persist below firstUpdateIndex \\
ConcreteProtocolStep & 14 & CPS2: A5 bridge (concrete$\to$abstract) \\
ElectionReachability & 12 & ER1--ER10: high-water-mark + shared-source \\
ElectionConcreteModel &  8 & ECM1--ECM7: hae $\to$ hqc\_preserved \\
AEBroadcastInvariant & 10 & ABI8: hqc\_preserved\_of\_broadcast \\
ElectionBroadcastChain &  6 & EBC4: broadcast seq $\Rightarrow$ HaeCovered \\
ElectionLifecycle &  7 & EL1--EL7: election epoch $\Rightarrow$ RaftReachable \\
MultiStepReachability &  7 & MS2: N-step end-to-end safety certificate \\
MaybePersistFUI   &  8 & FU1: snap blocks advance; FU8: fui $>$ persisted \\
RaftLogAppend     & 16 & RA1--RA9, P4--P7: append spec \\
ReadOnly          & 14 & RO8: advance removes ctx; RO14--RO15: inv \\
FindConflictByTerm & 10 & FCB3: maximality; FCB6: always some \\
ProgressTracker   & 26 & PT1--PT26: all\_wf invariant preservation \& integration \\
UncommittedState  & 18 & US16--US18: budget accept/reject characterisation \\
ConfigurationInvariants & 11 & CI5: voters$\cap$learners=$\emptyset$; CI11: strict vs.\ relaxed \\
UnstablePersistBridge &  8 & Bridge: MaybePersist$\leftrightarrow$LogUnstable \\
FindConflictCorrespondence &  2 & 20 \texttt{\#guard}, exact \\
MaybeAppendCorrespondence &  0 & 32 \texttt{\#guard}, exact \\
IsUpToDateCorrespondence  &  0 & 12 \texttt{\#guard}, exact \\
CommittedIndexCorrespondence & 0 & 11 \texttt{\#guard}, abstraction \\
LimitSizeCorrespondence   &  0 & 10 \texttt{\#guard}, abstraction \\
ConfigValidateCorrespondence & 0 & 12 \texttt{\#guard}, exact \\
InflightsCorrespondence   &  0 & 12 \texttt{\#guard}, abstraction \\
LogUnstableCorrespondence &  0 & 12 \texttt{\#guard}, exact \\
TallyVotesCorrespondence  &  0 & 10 \texttt{\#guard}, exact \\
VoteResultCorrespondence  &  0 & 15 \texttt{\#guard}, exact \\
HasQuorumCorrespondence   &  0 & 15 \texttt{\#guard}, exact \\
ReadOnlyCorrespondence    &  0 & 15 \texttt{\#guard}, exact \\
FindConflictByTermCorrespondence & 0 & 16 \texttt{\#guard}, abstraction \\
ProgressCorrespondence    &  0 & 51 \texttt{\#guard}, abstraction \\
ProgressTrackerCorrespondence & 0 & 44 \texttt{\#guard}, abstraction \\
MaybePersistCorrespondence &  0 & 21 \texttt{\#guard}, abstraction \\
MaybeCommitCorrespondence  &  0 & 17 \texttt{\#guard}, abstraction \\
RaftLogAppendCorrespondence & 0 & 21 \texttt{\#guard}, abstraction \\
MaybePersistFUICorrespondence & 0 & 25 \texttt{\#guard}, abstraction \\
ElectionCorrespondence    &  0 & 28 \texttt{\#guard}, abstraction \\
ConfigurationInvariantsCorrespondence & 0 & 18 \texttt{\#guard}, abstraction \\
JointVoteCorrespondence   &  0 & 18 \texttt{\#guard}, abstraction \\
QuorumRecentlyActiveCorrespondence &  0 & 14 \texttt{\#guard}, abstraction \\
HasNextEntriesCorrespondence & 0 & 33 \texttt{\#guard}, abstraction \\
NextEntriesCorrespondence &  0 & 28 \texttt{\#guard}, abstraction \\
UncommittedStateCorrespondence & 0 & 13 \texttt{\#guard}, abstraction \\
ProgressSetCorrespondence & 0 & 28 \texttt{\#guard}, abstraction \\
MemStorageCorrespondence  & 0 & 23 \texttt{\#guard}, abstraction \\
BroadcastLifecycle &  3 & BL1--BL3: BroadcastSeq$\to$RaftReachable bridge \\
HasNextEntries    & 14 & HNE1--HNE14: hasNextEntriesSince boolean spec \\
NextEntries       &  7 & NE1--NE7: nextEntriesSince slice model \\
ProgressSet       &   8 & PS1--PS8: composed progress set \\
MemStorage        &  14 & Compact/append contiguity \\
IsContinuousEnts  &   8 & ICE1--ICE8: batch concat validity \\
Restore           &   8 & RST1--RST8: snapshot restore spec \\
\midrule
\textbf{Total}    & \textbf{718} & \textbf{0 sorry; 625 \texttt{\#guard}} \\
\bottomrule
\end{tabular}
\caption{Full file inventory (all at Phase~5, 0 sorry). Correspondence files contribute \texttt{\#guard} assertions but no named theorems (except \texttt{FindConflictCorrespondence}, 2 theorems).}
\label{tab:inventory}
\end{table*}

\subsection{Bugs and Findings}
\label{sec:findings}

\paragraph{No implementation bugs found.}
All 718~theorems are consistent with the Rust implementation, as further
confirmed by 625 compile-time \texttt{\#guard} correspondence assertions
and 18+ Rust test functions.

\paragraph{Formulation bug caught and corrected.}
Early versions of \texttt{log\_matching\allowbreak\_property} (RSS3) and
\texttt{raft\_committed\allowbreak\_no\_rollback} (RSS4) claimed properties for
\emph{arbitrary} log states---with no hypotheses about how the logs were
constructed.  These statements are provably false (a log with random
contents satisfies no log-matching property).  The \texttt{sorry} mechanism
acted as a ``needs review'' marker: since \texttt{sorry} was required but
no proof could be found, the error was caught before it entered the proof
base.

The corrected versions carry explicit hypotheses:
\texttt{LogMatching\-InvariantFor} and
\texttt{NoRollback\-InvariantFor},
which are exactly the invariants preserved by valid Raft transitions.
Both were subsequently proved (run r130).

%% ============================================================
\section{Discussion}
%% ============================================================

\subsection{Proof Architecture and the Closed Protocol Gap}
\label{sec:reachability}

The \texttt{RaftReachable} model (Figure~\ref{fig:rr}) provides
honest accounting of what the safety proof depends on.  The five \texttt{step}
hypotheses (\texttt{hlogs'}, \texttt{hno\_overwrite}, \texttt{hqc\_preserved},
\texttt{hcommitted\_mono}, \texttt{hnew\_cert}) precisely state what a
protocol transition must satisfy to preserve \texttt{CommitCertInvariant}.

As of Run~126, \textbf{all five hypotheses have been fully discharged}
from concrete proved theorems:

\begin{itemize}\setlength\itemsep{0pt}
\item \texttt{hlogs'}: discharged by \texttt{ValidAEStep}, which models
  single-voter AppendEntries (CPS8/CPS9).
\item \texttt{hno\_overwrite}: discharged by CPS1
  (\texttt{validAEStep\_hno\_over\-write}), via the committed-$\leq$-prev
  precondition and CT2.
\item \texttt{hqc\_preserved}: discharged by EL1
  (\texttt{electionEpoch\_hqc\_pre\-served}), which chains election safety
  (RE5), the broadcast induction (EBC6), and the election lifecycle
  bridge (A7).
\item \texttt{hcommitted\_mono}: discharged by CPS11 from
  \texttt{ValidAEStep}.
\item \texttt{hnew\_cert}: discharged by CR8 (\texttt{CommitRuleValid})
  together with MC4 (A6 term safety: \texttt{maybeCommit} only commits
  entries from the current leader's term).
\end{itemize}

The key theorem is \texttt{fullProtocolStep\_safe} (EL7): given any
\texttt{RaftReachable} state, an \texttt{ElectionEpoch} (election +
broadcast), and a subsequent \texttt{ValidAEStep}, the resulting state
is cluster-safe.  No abstract axioms remain.

The proof chain for the hardest hypothesis (\texttt{hqc\_preserved})
proceeds: election safety (RE5, at most one leader per term) $\to$
vote-granting implies \texttt{isUpToDate} (RE7) $\to$
\texttt{BroadcastSeq} induction (EBC4) $\to$
\texttt{hqc\_preserved} from broadcast (EBC6) $\to$
election lifecycle bridge (EL1--EL7).

Concretely, \texttt{hno\_overwrite} corresponds to the Rust panic guard:
\begin{lstlisting}[language=Rust2]
if conflict != 0
   && conflict <= self.committed {
  panic!("entry {} conflict
    with committed log", conflict);
}
\end{lstlisting}
This panic is proved to be unreachable in a correct Raft execution by
\texttt{hno\_overwrite}.  Theorem RP8 (\texttt{raft\_transitions\_no\_rollback})
proves that \texttt{maybe\_append} satisfies \texttt{hno\_overwrite} given
the appropriate preconditions.

\subsection{Proof Utility Assessment}

The 718 theorems divide into two groups by structural role.
\textbf{The safety chain} (25 files, 341 theorems) forms a connected
dependency graph from quorum arithmetic through to
\texttt{fullProtocolStep\_safe} (EL7), the end-to-end safety result.
Every theorem in this chain has a genuine transitive dependency on the
lower layers; removing any would break the final result.
\textbf{Standalone module verifications} (25 files, 375 theorems) prove
properties of individual Raft subsystems---inflights, log-unstable,
progress tracking, read-only, flow control, etc.---that do not feed
into the safety chain but independently validate critical data-structure
invariants.  Both groups contribute to overall confidence.

In particular:

\begin{itemize}\setlength\itemsep{0pt}
\item HQ20 (\texttt{quorum\_intersection\_mem}) is the \emph{only} mathematical
  fact about quorums used in the RSS1 proof.  If it were wrong, RSS1 would
  fail.  Its proof is non-trivial (inclusion-exclusion + arithmetic), giving
  meaningful assurance.
\item MA13/MA14 (\texttt{MaybeAppend}) give a complete behavioural
  characterisation of AppendEntries log modification---the most frequent
  Raft protocol operation.  A bug in \texttt{maybe\_append} would cause
  one of these to fail.
\item SC4 (\texttt{SafetyComposition}) is the first theorem that cannot
  be stated within a single module; it requires both quorum arithmetic
  and committed-index computation (both Layer~2) to interact.
\item The standalone modules, while not in the safety chain, verify the
  data-structure invariants that the Raft implementation relies on in
  practice.  For example, the 40 \texttt{Inflights} theorems prove that
  the ring-buffer implementation computes identical results to a simple
  list model---a property that matters for operational correctness even
  though the safety proof abstracts over it.
\end{itemize}

The main remaining abstraction is the Rust implementation itself:
the Lean models capture the algorithms faithfully (as confirmed by
625 \texttt{\#guard} assertions), but do not use mechanical extraction
from Rust source.  Fully unconditional safety with respect to the
Rust binary would require mechanical extraction (e.g.\ via
Aeneas~\cite{aeneas}) or a verified compiler toolchain.

\subsection{Automation and Human Intervention}
\label{sec:interventions}

The entire FV effort was conducted by the \emph{Lean Squad} automated
agent, an LLM-driven agent that runs as a GitHub Actions workflow
(\S\ref{sec:workflow}).
Each run: reads persistent memory, draws two tasks from a phase-weighted
schedule, installs the Lean toolchain, executes the tasks (writing or
modifying Lean files, specs, or documentation), runs \texttt{lake build},
and creates a pull request with verified code.  A final status-issue
update always runs.

The human author's role was deliberately hands-off: he designed the
Lean Squad workflow itself (\S\ref{sec:workflow}), installed it on the
repository, and reviewed merged PRs.  Over the entire 129-run project,
his total intervention consisted of four brief GitHub issue comments:

\paragraph{Intervention 1 (Issue \#94 comment, run~$\approx$28).}
After 370 theorems were proved but no state-machine-level safety
theorem existed, the human wrote:
\begin{quote}\small
\texttt{/lean-squad Address these concerns from critique.md,
plan out rest of work, write theorems even if unproven,
get the project on Trajectory to completion}
\end{quote}
This triggered \texttt{RaftSafety.lean} and
\texttt{CrossModuleComposition.lean} (run~r128), adding 25 new
theorems including the first end-to-end cluster safety results.

\paragraph{Intervention 2 (Issue \#144, run~$\approx$29).}
After the initial proof effort was declared complete, the human wrote:
\begin{quote}\small
\texttt{/lean-squad Now FV is complete, please update critique,
correspondence etc.\ and write a final REPORT.md project report.
Include mermaid diagrams}
\end{quote}
This produced the project report and updated documentation artefacts.

\paragraph{Intervention 3 (Issue \#149, run~$\approx$35).}
The most substantive intervention: the human pasted a detailed
external critique identifying the structural gap in the
\texttt{RaftReachable.step} hypotheses (specifically that
\texttt{hqc\_preserved} required leader completeness, which
required an election model not yet formalised) and wrote:
\begin{quote}\small
\texttt{/lean-squad Continue the work, adjusting report, targets
and replanning all future work, based on this critique}
\end{quote}
This triggered the agent to incorporate the gap analysis and
replan, eventually producing the election model (Layer~7),
concrete protocol bridge, and leader completeness proofs.

\paragraph{Intervention 4 (Issue \#255, run~$\approx$125).}
The human asked the agent to assess composition completeness
and then wrote:
\begin{quote}\small
\texttt{/lean-squad OK go through and make progress on whatever
else needs doing to close out the work}
\end{quote}
This directly produced \texttt{ElectionLifecycle.lean}
(EL1--EL7), closing the A7 gap and fully discharging all
five \texttt{RaftReachable.step} hypotheses.

\paragraph{What is remarkable} is not that four interventions occurred,
but what did \emph{not} require human intervention: the choice of all
26 FV targets, the layered proof architecture, all 25 informal
specifications, all Lean formalisations, all 718 proofs, the 28
correspondence-validation files, the proof utility critique, the project
report, and even an initial draft of this paper were produced autonomously.
The human never selected a target, wrote a specification, suggested a
proof strategy, or debugged a failing proof.

Across 129 runs spanning $\approx$18 weeks, the agent produced:
\begin{itemize}\setlength\itemsep{0pt}
  \item 25 informal specification documents
  \item 79 Lean~4 source files ($\approx$21{,}000 lines)
  \item 718 machine-checked theorems
  \item 625 compile-time \texttt{\#guard} correspondence assertions
  \item 18+ Rust correspondence test functions
  \item 1 formulation bug caught and corrected
  \item 0 implementation bugs
  \item All 5 \texttt{RaftReachable.step} hypotheses fully discharged
\end{itemize}

The key insight enabling this level of autonomy is the workflow's
phase-weighted task selection (\S\ref{sec:tasksel}): the agent always
has a well-defined next action, and persistent memory ensures each run
builds on a growing base of named lemmas and established
infrastructure.  The layered architecture emerged organically from the
agent's interaction with its own memory and the growing codebase.

\subsection{Lean Stdlib vs.\ Mathlib}

We deliberately chose Lean~4 stdlib only, without Mathlib.  This keeps
the proof artefact self-contained and avoids version-pinning issues
with Mathlib's rapid development.  The cost is the absence of mature
list-processing lemmas and richer automation.

In practice, most proofs required only \texttt{omega}, \texttt{simp},
\texttt{by\_cases}, and structural induction---all available in stdlib.
The only Lean~4 stdlib pitfalls encountered were: \texttt{lemma} not
being a top-level keyword (use \texttt{theorem} instead),
\texttt{split\_ifs} and \texttt{push\_neg} not being available
(use \texttt{by\_cases} and \texttt{omega}), and \texttt{le\_rfl}
not existing (use \texttt{Nat.le\_refl \_}).

\subsection{Lessons Learned}

\paragraph{Phase-weighted task selection enables autonomy.}
The workflow's ten-task structure with adaptive weights means the
agent always has a well-defined next action.  The agent does not need
a human to tell it what to do next---the weight computation, combined
with persistent memory, provides sufficient direction.  The layered
proof architecture emerged organically from this process rather than
being prescribed.  The correspondence-validation layer (Layer~8) emerged
when the agent recognised the need to ground abstract Lean models in
concrete Rust behavior, and the election-lifecycle layer emerged
when the agent set about closing the \texttt{RaftReachable.step}
hypotheses.

\paragraph{Persistent memory is essential for cumulative progress.}
An LLM agent can prove theorems in a new module if all dependencies
are already available as named lemmas.  Persistent memory ensures
the agent knows which lemmas exist, which targets are complete, and
what strategies have worked or failed in prior runs.

\paragraph{The \texttt{sorry} discipline catches formulation bugs.}
Proving with \texttt{sorry} as a placeholder and then trying to remove it
is a reliable method for finding specification errors.  Both RSS3 and RSS4
were caught this way.

\paragraph{One module per run, not the whole codebase.}
Each run focused on a single target.  This made the Lean context small
enough for the LLM to reason about correctly, and produced clean,
reviewable pull requests.

\paragraph{Formal models must be honest about what they omit.}
The CORRESPONDENCE.md document records all known divergences between
Lean models and Rust source.  Making these explicit---rather than
claiming full correspondence---is essential for correctly assessing
proof value.

\paragraph{Negative results have value.}
Finding no implementation bugs in 718~theorem-checks is itself evidence
of implementation quality.  The absence of counterexamples in the
verified paths raises confidence, even without a fully concrete model.

\paragraph{Not all theorems feed the top-level result, and that is fine.}
Only 25 of the 50 proof files (341 of 716 proof-file theorems) form the
dependency chain to \texttt{fullProtocolStep\_safe}.  The remaining 25
files (375 theorems) verify standalone data-structure invariants.  An
LLM agent exploring a codebase will naturally verify what it encounters,
producing useful module-level proofs even when they do not connect to
the global safety theorem.  This mirrors human FV practice:
not every invariant proved in seL4 feeds the top-level kernel safety
property.  The standalone proofs still have value as independently
verified specifications.

\subsection{Cost}
\label{sec:cost}

Each of the 129 Lean Squad runs consumed approximately \$3 of LLM
token cost at Claude Sonnet~4.6 pricing (at time of publication), for a total of
approximately~\$400.  This covers the full pipeline: reading
persistent memory, researching the codebase, selecting tasks,
installing the Lean toolchain, writing or modifying Lean files,
running \texttt{lake build}, and creating pull requests.  The
GitHub Actions compute time (Linux runners) is negligible relative
to token cost.

For comparison, IronFleet~\cite{ironfleet} required person-years of
expert proof engineering in Dafny; Woos~et~al.~\cite{vard} reported
$\sim$50{,}000 lines of Coq proof.  Even valuing human proof-engineer
time conservatively, the dollar cost of comparable manual efforts
exceeds the LLM cost by orders of magnitude.  The trade-off is clear:
LLM-driven FV substitutes cheap, parallelisable compute for scarce,
expensive human expertise.  As model prices continue to fall, the
economics will only improve.

\subsection{Training Data Overlap}
\label{sec:training}

It is virtually certain that Claude Sonnet~4.6 was trained on published
papers, source code, and proof artefacts related to the formal
verification of Raft---including Ongaro's dissertation~\cite{ongaro2014dissertation},
the Verdi/VRaft Coq development~\cite{vard}, TLA+ specifications, and
the \texttt{raft-rs} source code itself.  The degree to which this
prior exposure influenced the agent's target selection, proof
architecture, and tactic choices cannot easily be measured.

We note two mitigating factors.  First, the proofs are machine-checked
by Lean's kernel: regardless of how the agent arrived at a proof term,
its correctness is independently verified.  Training-data influence
cannot produce a false positive.  Second, several architectural
decisions---stdlib-only (no Mathlib), the layered decomposition, the
\texttt{RaftReachable} abstract model---have no direct precedent in the
prior Raft verification literature, suggesting that the agent did more
than reproduce memorised material.  Nevertheless, this overlap should
be borne in mind when assessing the generalisability of our results to
domains with less existing formal-verification literature.

%% ============================================================
\section{Related Work}
\label{sec:related}
%% ============================================================

We situate our contribution with respect to six threads of prior work.

\subsection{Formal Verification of Raft}

The closest predecessor is Woos~et~al.~\cite{vard}, who gave the first
machine-checked proof of state-machine safety for Raft in Coq using the
Verdi framework~\cite{verdi}.  Their proof discovered 90 system invariants
and extracted a runnable OCaml implementation.  A key methodological
contribution was \emph{planning for change}: factoring proofs so that
strengthening one invariant does not cascade failures throughout the
development.  Our work differs in three respects: (i)~we target a
\emph{production} Rust implementation (\texttt{raft-rs}) rather than an
extracted model; (ii)~we use Lean~4 with stdlib only; and (iii)~the
entire effort is LLM-automated.  Our layered architecture echoes
their information-hiding methodology, but arises from the need to give
the LLM agent a well-defined proof frontier rather than from human
proof-maintenance concerns.

Subsequent Raft verification efforts have used diverse formalisms.
Evrard~\cite{evrard2020raft} modelled Raft in LNT for model-checking
with CADP.  Bao~et~al.~\cite{bao2022raft} and
Ishibashi~and~Ogata~\cite{ishibashi2023raft} model-checked Raft leader
election in NuSMV and Maude respectively.
Sanabria~et~al.~\cite{sanabria2024raft} verified Raft leader election
safety using the Athena proof system.
Bora~et~al.~\cite{bora2024raft} modelled Raft in mCRL2, enabling
behavioural equivalence checking.
Xiong~et~al.~\cite{xiong2025recraft} formalised split, merge, and
membership changes for Raft (ReCraft), complementing our focus on
log safety.  None of these target a production implementation or
produce end-to-end safety proofs at the level of
\texttt{raftReachable\_safe}.

Ongaro's TLA+ specification~\cite{ongaro2014dissertation} provides
a complementary model-checking approach.  TLA+
can exhaustively check finite instances but does not produce
machine-checked proofs over unbounded state spaces.  Our Lean proofs
hold for arbitrary cluster sizes and unbounded execution traces,
within the \texttt{RaftReachable} model.

\subsection{Verified Distributed Systems}

IronFleet~\cite{ironfleet} verified Paxos-based replication and a
distributed key-value store in Dafny with TLA-style temporal
reasoning.  The Disel framework~\cite{sergey2018disel} introduced
dependent separation protocols for modular distributed-system
verification in Coq.  Sprenger~et~al.~\cite{sprenger2020igloo}
proposed Igloo, combining compositional refinement with separation
logic for distributed-system verification.
P\^irlea~and~Sergey~\cite{pirlea2018blockchain} mechanised blockchain
consensus in Coq, demonstrating that protocol-level safety proofs
are tractable in proof assistants.

More recently, Honor\'e~et~al.\ developed a series of verified
distributed-object frameworks: Adore~\cite{honore2022adore} for
certified reconfiguration, ADO~\cite{honore2021ado} for
fault-aware compositional verification, and
AdoB~\cite{honore2024adob} for bridging benign and Byzantine
consensus.  Qiu~et~al.\ verified linearisable Byzantine distributed
objects (LiDO~\cite{qiu2024lido}) and extended this to DAG-based
consensus (LiDO-DAG~\cite{qiu2025lidodag}).
Kim~et~al.~\cite{kim2025verirt} presented VeriRT, an end-to-end
verification framework for real-time distributed systems.
Lewchenko~et~al.~\cite{lewchenko2025bolton} verified bolt-on strong
consistency specifications.
Zhao~et~al.~\cite{zhao2024byzantine} developed compositional
verification of composite Byzantine protocols.
Zhang~et~al.~\cite{zhang2024invariants} studied inductive invariants
for distributed protocols with a focus on proof ergonomics.

Two recent systems-level efforts are particularly relevant.
AutoMan~\cite{zhang2025automan} combines automatic code generation
with manual optimisations for verified distributed-system development,
achieving strong guarantees on production-grade systems.
Veil~\cite{pirlea2025veil} provides a framework for automated and
interactive verification of transition systems, offering an
alternative path to the reachability-style reasoning in our Layer~6.

Blockchain consensus verification has also grown rapidly:
Thomsen~and~Spitters~\cite{thomsen2021nakamoto} formalised
Nakamoto-style proof of stake in Coq;
Rashid~et~al.~\cite{rashid2025beacon} verified safety properties of
Ethereum~2.0 Beacon Chain epoch processing;
Schultz~et~al.~\cite{schultz2022reconfig} verified a distributed
dynamic reconfiguration protocol.

\subsection{Formal Verification of Rust}

Verifying Rust programs is an active area with several complementary
approaches.  RustBelt~\cite{jung2018rustbelt} established the
semantic foundations of Rust's type system safety guarantees in Coq
using Iris separation logic.
Prusti~\cite{astrauskas2022prusti} provides automated deductive
verification for Rust via Viper, targeting ownership and memory
safety.
Creusot~\cite{denis2022creusot} translates Rust to WhyML for
verification with Why3, handling borrows and ownership.
Verus~\cite{lattuada2023verus} embeds verification conditions
directly in Rust using linear ghost types and SMT solving.

Our approach is orthogonal: rather than verifying the Rust code
directly, we build a pure functional Lean model of the relevant
algorithms and verify properties of that model, documenting the
correspondence to Rust in \texttt{CORRESPONDENCE.md}.  This is
a deliberate trade-off: we lose mechanical connection to the Rust
source (which Aeneas~\cite{aeneas} could provide) but gain the full
power of Lean's dependent types and tactic framework for
protocol-level reasoning that SMT-based tools struggle with.

B\'il\'y~et~al.~\cite{bily2025rust} proposed a refinement methodology
for distributed Rust programs that could, in principle, bridge our
abstract Lean models with the concrete \texttt{raft-rs} code.
Combining their refinement approach with our verified safety
properties is a promising direction for future work.

\subsection{LLM-Based Theorem Proving}

Recent years have seen rapid progress in applying LLMs to formal
theorem proving.  LeanDojo~\cite{yang2024leandojo} introduced a
benchmark and retrieval-augmented approach for Lean proving, where
an LLM retrieves relevant premises from Mathlib before generating
tactic steps.  LEGO-Prover~\cite{wang2024legoprover} grows a
reusable library of verified lemmas across proof attempts,
conceptually similar to our persistent memory mechanism.
Baldur~\cite{first2023baldur} demonstrated whole-proof generation
and repair with LLMs for Isabelle/HOL.
Draft-Sketch-Prove~\cite{jiang2023dsp} uses informal proofs as
sketches that guide formal proof search.
Copra~\cite{thakur2024copra} employs an LLM agent with
backtracking for Lean theorem proving.
Polu~and~Sutskever~\cite{polu2022formal} showed that curriculum
learning over formal mathematics enables LLMs to solve increasingly
difficult problems.
Song~et~al.~\cite{song2024towards} explored LLMs as interactive
copilots for Lean users.

Bayaz\i t~et~al.~\cite{bayazit2025llm} conducted a case study on
the effectiveness of LLMs for verification with proof assistants,
finding that LLMs can discharge simple lemmas but struggle with
multi-step proofs requiring creative insight.  Our experience
\emph{diverges}: the Lean Squad agent completed all 718~theorems
including multi-step compositional proofs (Layers~4--7) that require
chaining lemmas across modules.  We attribute this difference to
four factors: (i)~the workflow's phase-weighted task selection
(\S\ref{sec:workflow}) decomposes the problem into tractable
increments, each with a well-defined interface to the next;
(ii)~persistent memory across runs enables cumulative progress
where each run builds on a growing base of named lemmas and
established infrastructure;
(iii)~the remarkable improvement in raw LLM reasoning power in the
intervening period---frontier models in 2025--2026 are substantially
more capable at multi-step logical reasoning than those available
during the Bayaz\i t~et~al.\ study; and
(iv)~the planning and tool-use capabilities of modern coding agents
(Copilot CLI, Codex CLI), which provide structured
multi-step reasoning, file navigation, and iterative error correction
that go well beyond single-prompt LLM interactions.

These systems differ from our approach in a fundamental way:
they target \emph{benchmark-style} theorem proving, where the goal
is to close a given lemma statement in an existing formalisation.
The human provides the theorem statement and the library context;
the LLM provides the proof.  In contrast, \emph{Lean Squad operates
at the project level}: it selects what to verify, writes the
specifications, builds the library from scratch, and proves the
theorems---all within a production codebase that has no prior
formalisation.

The closest analogue to our persistent-memory mechanism is
LEGO-Prover's growing library~\cite{wang2024legoprover}, but
LEGO-Prover operates within a single proof session, whereas Lean
Squad accumulates knowledge across 129 independent GitHub Actions runs
spanning weeks.  AlphaProof~\cite{trinh2024alphaproof} demonstrated
that reinforcement learning over formal proofs can achieve
competition-level mathematical reasoning, but targets Olympiad
geometry rather than software verification, and requires massive
compute for training.

Yao~et~al.~\cite{yao2024liveness} demonstrated mostly-automated
verification of liveness properties using ranking functions, and
Gopinathan~et~al.~\cite{gopinathan2023repair} developed
mostly-automated proof repair for verified libraries---both pointing
toward increasing automation in formal verification.
D\'iaz~et~al.~\cite{diaz2025incremental} introduced incremental
certified programming, which shares our goal of making cumulative
progress on large proof developments.

\subsection{Autonomous AI Agents for Software Engineering}

SWE-bench~\cite{jimenez2024swebench} established a benchmark for
evaluating LLM agents on real-world GitHub issues.
SWE-agent~\cite{yang2024sweagent} demonstrated that purpose-built
agent-computer interfaces improve automated issue resolution.
Agentless~\cite{xia2024agentless} showed that simpler, non-agentic
approaches can also resolve a significant fraction of issues.

These agents operate in a \emph{reactive} mode: given a bug report or
feature request, they produce a patch.  Lean Squad operates in a
\emph{proactive, scheduled} mode: it runs every 8 hours regardless
of external triggers, inspects the repository state, and advances
the project.  This distinction is important because formal
verification is not naturally triggered by external events---it is
a long-running engineering effort that requires sustained, cumulative
progress.

The phase-weighted task selection mechanism (\S\ref{sec:tasksel})
can be viewed as a form of \emph{automated project management}:
the workflow maintains a notion of project phase, allocates effort
accordingly, and adapts as the project matures.  This is, to our
knowledge, novel among LLM-agent architectures.

\subsection{Proof Engineering and Large-Scale Verification}

Ringer~et~al.~\cite{ringer2019qed} surveyed the landscape of
formally verified software engineering, identifying proof maintenance
as a central challenge---the same challenge that motivated the
``planning for change'' methodology of Woos~et~al.~\cite{vard}.
Dejon~et~al.~\cite{dejon2024codetoqed} provided a systematic guide
to proof engineering, and Luan~et~al.~\cite{luan2025compatibility}
studied compatibility issues across theorem prover versions, a
practical concern we sidestep by using Lean~4 stdlib only.
The proof-repair line of work by
Ringer~et~al.~\cite{ringer2021repair} and
Viola~et~al.~\cite{viola2025repair} addresses proof maintenance
across type equivalences, complementing our approach of isolating
change behind layer boundaries.

Fiat-Crypto~\cite{fiat-crypto} demonstrated automatic synthesis of
verified cryptographic arithmetic.  The Aeneas
toolchain~\cite{aeneas,charon} enables mechanical Lean extraction
from safe Rust; we note its applicability but do not use it in this
work.

To calibrate expectations, we compare the scale of our verification
with landmark projects.
seL4~\cite{klein2009sel4} verified a production OS microkernel: 8,700
lines of C, $\sim$200,000 lines of Isabelle/HOL proof, 11 person-years.
CertiKOS~\cite{gu2016certikos} verified a concurrent OS kernel with
6,500 lines of C and $\sim$100,000 lines of Coq.
CompCert~\cite{compcert} verified an optimising C compiler over
$\sim$100,000 lines of Coq.
IronFleet~\cite{ironfleet} verified Paxos replication with
$\sim$42,000 lines of Dafny proof.
Woos~et~al.~\cite{vard} verified simplified Raft with $\sim$50,000
lines of Coq across 90 invariants.

Our $\approx$21{,}000 lines of Lean and 718 theorems are modest by comparison.
However, the relevant metric is not raw scale but \emph{effort}:
those projects required person-years of expert proof engineering,
while ours required zero human proof engineering and approximately
129 LLM agent runs (each $\sim$15--30 minutes of wall-clock time)
plus two one-line human comments.  If the approach scales---and the
successful progression from Layer~1 through Layer~20 suggests it
can---the implications for the cost of formal verification are
significant.

\medskip\noindent
Our contribution is, to our knowledge, the first machine-checked
safety proof for a \emph{production} Raft implementation, using a
modern proof assistant without external libraries, and generated
entirely by an automated LLM agent.

%% ============================================================
\section{Conclusion}
%% ============================================================

Formal verification has historically required deep expertise in proof
assistants and months of dedicated human effort.  This paper demonstrates
that LLM-driven agents can eliminate that barrier.

Starting from zero, the Lean Squad agent---running as a scheduled
GitHub Actions job with no human in the loop---autonomously selected
26 verification targets in a production Raft implementation, designed
a layered proof architecture, wrote all specifications and proofs,
and produced 718 machine-checked theorems across 79 files in 129 runs,
culminating in the end-to-end safety theorem
\texttt{raftReachable\_safe} with zero \texttt{sorry} and all five
abstract \texttt{RaftReachable.step} hypotheses fully discharged from
concrete protocol theorems, plus the A5
bridge (CPS2) connecting concrete AppendEntries steps to the abstract
reachability model, and an additional 28 correspondence-validation
files providing 625 compile-time \texttt{\#guard} assertions and 18+
Rust test functions confirming that the Lean models agree with the
Rust implementation on concrete inputs.
The human author's total involvement was: designing the reusable
workflow, merging PRs, and writing two short comments.
He required no Lean expertise, no formal-methods training, and
invested near-zero hours of proof-engineering labour.

The key enablers are: (i)~a phase-weighted task selection mechanism
that gives the agent a well-defined next action at every stage of the
project; (ii)~persistent memory that accumulates proved lemmas,
target status, and proof strategies across independent runs; and
(iii)~the dramatically improved reasoning capabilities of modern
coding agents (Copilot CLI with Claude Sonnet~4.6 in this project).

The proof architecture is layered, honest, and extensible: each layer's
assumptions are explicit, and all five \texttt{RaftReachable.step}
hypotheses have been discharged from concrete protocol theorems.
The Lean Squad workflow is open-source
and installs on any GitHub repository with two commands
(\S\ref{sec:install}), requiring only an LLM API token.
We believe this shifts the economics of formal verification:
the cost moves from scarce human proof-engineering expertise to
commodity LLM compute, making machine-checked safety proofs a
practical option for ordinary software teams.

Future work on Lean Squad includes
reliably applying the Aeneas~\cite{aeneas} toolchain to mechanically
extract Lean models from safe Rust; generalising it to use a broader
range of formal verification tools and methodology; and applying Lean
Squad to other production codebases to test generality.

\begin{acks}
\emph{This paper was co-authored by Don Syme, LeanSquad (an autonomous
AI automation workflow using GitHub Agentic Workflows), Copilot CLI
(GitHub's AI coding agent), and Claude Sonnet~4.6 (Anthropic's LLM).  Don Syme designed the Lean Squad workflow
and made two brief interventions; all 26 target selections, all
specifications, all 718 proofs, the layered architecture, and the
correspondence-validation harness were produced autonomously by the AI
agent.  All theorems are machine-checked by Lean~4.30.0-rc2 with zero
\texttt{sorry}, and all five \texttt{RaftReachable.step} hypotheses
have been fully discharged.  The proof artefacts, agent configuration, and full run
history are available in the repository \texttt{dsyme/raft-lean-squad}.
The Lean Squad workflow is
open-source at \texttt{githubnext/\allowbreak agentics/\allowbreak workflows/\allowbreak lean-squad.md}.}
\end{acks}

\bibliographystyle{ACM-Reference-Format}
\bibliography{paper}

\end{document}
